% =====================================================================
%  Thème 3 — Colinéarité et parallélisme — Niveau FACILE — Exercices
% =====================================================================
\documentclass[11pt,a4paper]{article}
\usepackage[utf8]{inputenc}
\usepackage[T1]{fontenc}
\usepackage{lmodern}
\usepackage[french]{babel}
\usepackage{amsmath,amssymb,mathtools}
\usepackage{xcolor}
\usepackage{enumitem}
\usepackage{fancyhdr}
\usepackage[a4paper,margin=2.1cm,headheight=26pt,headsep=10pt]{geometry}

\definecolor{primary}{RGB}{222,70,80}
\definecolor{primarydark}{RGB}{150,30,42}

\pagestyle{fancy}\fancyhf{}
\fancyhead[L]{\small\textcolor{primarydark}{\textbf{Fiche 2} — Calcul vectoriel}}
\fancyhead[R]{\small\textcolor{primarydark}{\textsc{Thème 3 — Facile — Exercices}}}
\fancyfoot[C]{\small\thepage}
\renewcommand{\headrulewidth}{0.8pt}

\newcommand{\vc}[1]{\overrightarrow{#1}}
\providecommand{\norm}[1]{\left\lVert #1\right\rVert}
\newcommand{\exo}[1]{\medskip\noindent{\color{primary}\bfseries\large Exercice #1.}\par\nopagebreak\smallskip}

\begin{document}
\begin{center}
{\LARGE\bfseries\color{primarydark}Colinéarité et parallélisme — Niveau Facile}\\[2pt]
{\color{primary}\rule{0.5\linewidth}{1.1pt}}
\end{center}
\vspace{1mm}
{\small\itshape Objectif : reconnaître deux vecteurs colinéaires, appliquer le test du déterminant,
tester l'alignement et lire un vecteur directeur.}
\vspace{2mm}

\exo{1} \textbf{Colinéaires ou non (par un coefficient $k$) ?}
Dans chaque cas, dire si $\vec v=k\,\vec u$ pour un réel $k$, et donner $k$ le cas échéant.
\begin{enumerate}[label=\textbf{\alph*)},leftmargin=2em,itemsep=2pt]
  \item $\vec u\,(2;3)$ et $\vec v\,(6;9)$.
  \item $\vec u\,(1;-2)$ et $\vec v\,(-3;6)$.
  \item $\vec u\,(2;3)$ et $\vec v\,(4;5)$.
\end{enumerate}

\exo{2} \textbf{Test du déterminant $xy'-x'y$.}
Calculer le déterminant de chaque couple et conclure sur la colinéarité.
\begin{enumerate}[label=\textbf{\alph*)},leftmargin=2em,itemsep=2pt]
  \item $\vec u\,(3;4)$ et $\vec v\,(6;8)$.
  \item $\vec u\,(1;2)$ et $\vec v\,(3;5)$.
  \item $\vec u\,(-2;5)$ et $\vec v\,(4;-10)$.
\end{enumerate}

\exo{3} \textbf{Alignement de trois points.}
Dans chaque cas, dire si $A$, $B$, $C$ sont alignés (calculer $\vc{AB}$, $\vc{AC}$ et leur déterminant).
\begin{enumerate}[label=\textbf{\alph*)},leftmargin=2em,itemsep=2pt]
  \item $A\,(1;1)$, $B\,(3;2)$, $C\,(7;4)$.
  \item $A\,(0;0)$, $B\,(2;3)$, $C\,(4;5)$.
\end{enumerate}

\exo{4} \textbf{Vecteur directeur d'une droite.}
Donner un vecteur directeur de chaque droite (rappel : $ax+by+c=0$ donne $\vec u\,(-b;a)$).
\begin{enumerate}[label=\textbf{\alph*)},leftmargin=2em,itemsep=2pt]
  \item $d_1:\ 2x+3y-5=0$.
  \item $d_2:\ x-4y+1=0$.
  \item $d_3:\ y=2x+1$. \emph{(Réécrire d'abord sous la forme $ax+by+c=0$.)}
\end{enumerate}

\exo{5} \textbf{Coefficient directeur (pente).}
\begin{enumerate}[label=\textbf{\alph*)},leftmargin=2em,itemsep=2pt]
  \item Pente de la droite passant par $A\,(1;2)$ et $B\,(4;8)$.
  \item Pente d'une droite de vecteur directeur $\vec u\,(3;6)$.
  \item Pente de la droite d'équation $3x+2y-6=0$.
\end{enumerate}

\end{document}
